%! Absolute Value Functions & Transformations — Unit 1, Lesson 1.3 — Homework
% Gradual release: the "you do alone" block. Homework is generated per lesson (user direction,
% 2026-08-31); the teacher may substitute a DeltaMath set where the content is available there.
\documentclass[10pt]{article}
\usepackage{algebra2-article}
\usepackage{algebra2-boxes}
% Coordinate grid: {xmin}{xmax}{ymin}{ymax}
\newcommand{\lgrid}[4]{%
  \draw[step=1, linegray!40, very thin] (#1,#3) grid (#2,#4);
  \draw[->, semithick, charcoal] (#1,0)--(#2,0) node[below right, font=\scriptsize]{$x$};
  \draw[->, semithick, charcoal] (0,#3)--(0,#4) node[left, font=\scriptsize]{$y$};}

\begin{document}

\pageheader{Unit 1, Lesson 1.3}{Homework}

\vspace{0.10in}

\begin{notesbox}{Practice}
Show your work. For every rule, remember: in $a|x-h|+k$ the vertex is $(h,k)$ --- set the inside
to zero to find $h$ --- the axis is $x=h$, the sign of $a$ says up or down, and $|a|$ says narrow
or wide.

\begin{enumerate}[leftmargin=1.9em, itemsep=10pt]
  \item \textbf{Read the rule.} Fill in every feature.
    \begin{center}
    \renewcommand{\arraystretch}{1.4}
    \begin{tabularx}{\linewidth}{@{}l X X X X X@{}}
    & \textbf{Vertex} & \textbf{Axis} & \textbf{Opens} & \textbf{Max/min value} & \textbf{Range} \\
    (a) $y=|x-4|+2$ & \blank{1.3cm} & \blank{1.0cm} & \blank{1.0cm} & \blank{1.4cm} & \blank{1.3cm} \\
    (b) $y=-\tfrac12|x+1|+3$ & \blank{1.3cm} & \blank{1.0cm} & \blank{1.0cm} & \blank{1.4cm} & \blank{1.3cm} \\
    (c) $y=3|x|-4$ & \blank{1.3cm} & \blank{1.0cm} & \blank{1.0cm} & \blank{1.4cm} & \blank{1.3cm} \\
    \end{tabularx}
    \end{center}

  \boxguard[18]
  \item \textbf{One sign apart.} Compare the two rules, and notice what changes.
    \begin{center}
    \renewcommand{\arraystretch}{1.4}
    \begin{tabularx}{\linewidth}{@{}X|X@{}}
    \textbf{(a) $y=3|x-1|-2$} & \textbf{(b) $y=-3|x-1|-2$} \\
    \hline
    Vertex \blank{1.3cm}; opens \blank{1.1cm} & Vertex \blank{1.3cm}; opens \blank{1.1cm} \\
    $-2$ is a \blank{1.5cm}; range \blank{1.5cm} & $-2$ is a \blank{1.5cm}; range \blank{1.5cm} \\
    \end{tabularx}
    \end{center}
    (c) The two graphs share a vertex. Explain why that same point is a minimum for one and a
    maximum for the other. \writelines{2}

  \boxguard[20]
  \item \textbf{Match, then complete a table.}
    \begin{enumerate}[label=(\alph*), leftmargin=1.9em, itemsep=4pt]
      \item Write the letter of each rule's graph.
        \begin{center}
        \begin{tikzpicture}[scale=0.34]
          \begin{scope}
            \lgrid{-1}{5}{-1}{4}
            \draw[very thick, forest] (-1,3)--(2,0)--(5,3);
            \node[charcoal, font=\scriptsize] at (2,-2.4){\textbf{W}};
          \end{scope}
          \begin{scope}[xshift=8.5cm]
            \lgrid{-3}{3}{-3}{3}
            \draw[very thick, forest] (-3,1)--(0,-2)--(3,1);
            \node[charcoal, font=\scriptsize] at (0,-4.6){\textbf{X}};
          \end{scope}
          \begin{scope}[xshift=17cm]
            \lgrid{-3}{3}{-4}{2}
            \draw[very thick, forest] (-3,-3)--(0,0)--(3,-3);
            \node[charcoal, font=\scriptsize] at (0,-5.6){\textbf{Y}};
          \end{scope}
        \end{tikzpicture}
        \end{center}
        $y=|x-2|$ is \blank{0.9cm} \quad $y=|x|-2$ is \blank{0.9cm} \quad $y=-|x|$ is \blank{0.9cm}
      \item Complete the table for $y=|x-1|+2$, then read the vertex off it.
        \begin{center}
        \renewcommand{\arraystretch}{1.2}
        \begin{tabular}{|c|c|c|c|c|c|}
        \hline
        $x$ & $-1$ & $0$ & $1$ & $2$ & $3$ \\ \hline
        $y$ & \blank{0.7cm} & \blank{0.7cm} & \blank{0.7cm} & \blank{0.7cm} & \blank{0.7cm} \\ \hline
        \end{tabular}
        \end{center}
        Vertex \blank{1.4cm} --- the row where $y$ stops falling and starts rising.
    \end{enumerate}

  \boxguard[16]
  \item \textbf{The odd ones out.} Answer each from the picture, not a memorised rule.
    \begin{enumerate}[label=(\alph*), leftmargin=1.9em, itemsep=4pt]
      \item $y=|x|+k$ crosses the $x$-axis \blank{1.0cm} times when $k<0$, \blank{1.0cm} time(s)
            when $k=0$, and \blank{1.0cm} times when $k>0$.
      \item For which values of $a$ is $y=a|x|$ \emph{wider} than the parent? \blank{2.4cm}
            \; What happens to the graph when $a=0$? \blank{4.0cm}
    \end{enumerate}

  \boxguard[22]
  \item \textbf{Model it.} A drone flies in a straight line past a tower. Its distance from the
        tower, in metres, after $t$ seconds is $D(t)=15\,|t-6|$.
    \begin{enumerate}[label=(\alph*), leftmargin=1.9em, itemsep=4pt]
      \item The vertex is \blank{1.4cm}. In the story it means: \blank{6.0cm}
      \item When is the drone $45$ m from the tower? Solve $D(t)=45$.
        \begin{work}
          15\,|t-6| &= 45 \\
          |t-6| &= 3 \\
          t - 6 &= 3 \quad\text{or}\quad t - 6 = -3 \\
          t &= 9 \quad\text{or}\quad t = 3
        \end{work}
      \item If the drone flew \emph{twice as fast}, which number in the rule changes, and how does
            the V change? \writelines{2}
    \end{enumerate}

  \item \textbf{Multiple choice (SOL-style).} Which describes how the graph of $y=|x+4|-2$ is
        obtained from the graph of $y=|x|$?
    \begin{enumerate}[label=\Alph*., leftmargin=2.2em, itemsep=1pt]
      \item Shift right $4$ and down $2$
      \item Shift left $4$ and down $2$
      \item Shift left $4$ and up $2$
      \item Shift right $4$ and up $2$
    \end{enumerate}
    Say in one sentence how you ruled out one of the others. \writelines{2}
\end{enumerate}
\end{notesbox}

\boxguard[18]
\begin{extensionbox}
\textbf{Build and reason.}
\begin{enumerate}[label=(\alph*), leftmargin=1.9em, itemsep=6pt]
  \item A V opens \emph{up}, has vertex $(-2,1)$, and passes through $(0,5)$. Find $a$ and write
        the rule in vertex form.
        \begin{work}
          5 &= a\,|0-(-2)| + 1 \\
          5 &= 2a + 1 \\
          4 &= 2a \\
          a &= 2 \\
          y &= 2|x+2| + 1
        \end{work}
  \item For $y=a|x-h|+k$ with $a<0$, what is the range, in terms of $k$? Explain from the
        picture. \writelines{2}
\end{enumerate}
\end{extensionbox}

\begin{spiralbox}
\textbf{Coming next --- Lesson 1.4: Piecewise \& Step Functions.} Today's V is secretly
\emph{two} lines pasted at the vertex: $y=-x$ on the left and $y=x$ on the right. Next you will
write functions built from \emph{different rules on different pieces of the domain} --- the V is
the first one you already know --- and meet the greatest-integer (step) function, whose graph is
a staircase.
\end{spiralbox}

\end{document}
